\chapter{Alpha à court horizon}\label{ch:hf:short-horizon-alpha}

La file d'achat pèse trois fois la file de vente, et le future a monté il y a une demi-seconde. Un teneur de marché qui ignore les deux vendra au cours vendeur
à quelqu'un qui ne l'a pas fait~: dans le marché simulé de ce chapitre, où un fonds suit son future indiciel avec une demi-seconde de retard, cinq variables du carnet, les
transactions et le future prévoient le mid du fonds deux secondes à l'avance, avec un coefficient d'information de 0.48. C'est une borne supérieure~: le mid de ce
marché retarde sur son prix efficient pendant plusieurs secondes, ce que les carnets réels font moins souvent. Retirer la cotation que la prévision désigne
comme devant être prise pour cible fait passer le mark-out à deux secondes des exécutions du teneur de marché de 0.29 à 0.45 tick par action. La prévision perd à peine
de sa force en un message~; en cinquante, elle en perd plus d'un tiers.

\section{Ce qui prédit la seconde suivante}

\begin{definition}[Alpha à court horizon]\label{def:hf:short-horizon-alpha:alpha}
L'\emph{alpha à court horizon}\index{alpha à court horizon} est une prévision de la variation du prix d'un instrument sur les prochaines millisecondes à secondes, construite
à partir de l'état de son carnet d'ordres, de ses transactions récentes et des mouvements récents d'instruments liés, et utilisée pour décider comment coter et quand négocier.
\end{definition}

Il se distingue de l'alpha intrajournalier du livre~8, chapitre~14, par son horizon et son usage. À l'échelle des minutes, une prévision rapporte assez pour franchir la fourchette
et s'échange comme n'importe quel autre signal. À l'échelle des secondes, une prévision est rarement assez grande pour payer la fourchette~; son usage principal est
de décider lesquelles de ses propres cotations le teneur de marché laisse dans le carnet. Les prédicteurs sont les variables de carnet et de flux des chapitres~8 à~10 du livre~7~: le déséquilibre des files,
le déséquilibre du flux d'ordres, le flux signé des transactions, l'avance d'un instrument lié et la variation récente de l'instrument lui-même.

\section{Construire la prévision~: carnet, flux, instruments liés}

Le marché de ce chapitre comporte deux instruments reflétant un même prix efficient. A mène, le future indiciel~; B, le fonds, suit avec une demi-seconde de retard, chacun
avec son propre flux d'ordres. Le teneur de marché cote B et lit le carnet de A. Après chaque message de B, cinq variables sont calculées sur une fenêtre d'une seconde
(voir le \cref{lst:hf:short-horizon-alpha:features})~: le déséquilibre des files $(Q^b-Q^a)/(Q^b+Q^a)$~; le déséquilibre du flux d'ordres de Cont, Kukanov et Stoikov,
qui ont constaté que les variations de prix sur de courts intervalles sont linéaires en lui, avec une pente inversement proportionnelle à la profondeur~; le volume agressif net~; la variation du mid de A~;
et la variation du mid propre à B. Une régression ridge sur quatre sessions de calibration de vingt minutes prévoit la variation du mid de B sur les deux secondes suivantes.

Sur deux sessions de validation (coefficients d'information, signes tels qu'ajustés, prévision à deux secondes)~:

\begin{center}
\small
\setlength{\tabcolsep}{4pt}%
\begin{tabular}{@{}>{\raggedright\arraybackslash}p{3.0cm}>{\centering\arraybackslash}p{1.85cm}>{\centering\arraybackslash}p{1.85cm}>{\centering\arraybackslash}p{1.85cm}>{\centering\arraybackslash}p{1.85cm}>{\centering\arraybackslash}p{1.85cm}>{\centering\arraybackslash}p{1.85cm}@{}}
\toprule
& déséquilibre & déséq. flux d'ordres & flux de transactions & mouvement du leader & mouvement propre & combiné\\
\midrule
coefficient d'information & 0.320 & 0.285 & 0.077 & 0.288 & 0.150 & 0.480\\
par unité (ticks) & 0.30 & 0.0045 & 0.0057 & 0.34 & 0.27 & \\
\bottomrule
\end{tabular}
\end{center}

La deuxième ligne se lit ainsi~: un carnet entièrement déséquilibré vers l'achat prévoit $+0.30$ tick~; un tick de mouvement de A prévoit $+0.34$ tick de mouvement de B~; un lot
de déséquilibre du flux d'ordres, $+0.0045$. Le mouvement propre entre avec un signe positif~: dans ce marché, le mid suit une tendance pendant des secondes après qu'un prix efficient a sauté
(le livre~7 relève dans celui-ci une autocorrélation d'ordre un de 0.44 des rendements du mid sur deux secondes), et c'est cette tendance qui rend chaque coefficient ici plus grand que
dans les carnets réels. Il faut considérer les chiffres comme une borne supérieure et retenir le classement~: le carnet et le leader portent la prévision, le flux de transactions très peu.

\section{Horizon, décroissance et évaluation à haute fréquence}

Le coefficient d'information dépend de l'horizon (voir la \cref{fig:hf:short-horizon-alpha:ic})~: 0.28 à une demi-seconde, 0.48 à deux secondes, un maximum de 0.54 à
cinq secondes, 0.50 à dix secondes. Le déséquilibre culmine tôt puis s'estompe~; le déséquilibre du flux d'ordres continue de gagner jusqu'à dix secondes. Kolm, Turiel et Westray, en prévoyant 115
actions du Nasdaq à partir du flux d'ordres au moyen de réseaux de neurones, ont établi que l'horizon effectif de ces prévisions vaut environ deux variations moyennes de prix~; dans ce
marché, deux variations de prix prennent quelques secondes, et c'est à ce niveau que la courbe s'aplatit.

\begin{omfigure}
\begin{tikzpicture}
\begin{axis}[width=11cm,height=5.4cm,xlabel={\small horizon de prévision (secondes)},ylabel={\small coefficient d'information},xmode=log,log basis x=10,
  xmin=0.4,xmax=12,xtick={0.5,1,2,5,10},xticklabels={0.5,1,2,5,10},ymin=0,ymax=0.6,tick label style={font=\footnotesize},grid=major,
  grid style={gray!25},table/col sep=comma,legend style={font=\scriptsize,at={(0.5,-0.3)},anchor=north},legend columns=3]
\addplot[black,very thick,mark=*,mark size=1.3pt] table[x=h,y=combined]{figdata/hft/07-short-horizon-alpha/ic.csv};
\addplot[omDef,thick,mark=square*,mark size=1.1pt] table[x=h,y=imbalance]{figdata/hft/07-short-horizon-alpha/ic.csv};
\addplot[omThm,thick,dashed,mark=triangle*,mark size=1.3pt] table[x=h,y=ofi]{figdata/hft/07-short-horizon-alpha/ic.csv};
\addplot[omProp,thick,mark=diamond*,mark size=1.3pt] table[x=h,y=lead]{figdata/hft/07-short-horizon-alpha/ic.csv};
\addplot[gray,thick,dotted,mark=o,mark size=1.1pt] table[x=h,y=own]{figdata/hft/07-short-horizon-alpha/ic.csv};
\addplot[gray!60,thick,mark=x,mark size=1.3pt] table[x=h,y=flow]{figdata/hft/07-short-horizon-alpha/ic.csv};
\legend{combiné,déséquilibre,déséquilibre du flux d'ordres,mouvement du leader,mouvement propre,flux de transactions}
\end{axis}
\end{tikzpicture}

\omcaption{Coefficient d'information hors échantillon de chaque variable et de la combinaison ridge, prévoyant la variation du mid du fonds sur chaque
horizon (échelle logarithmique), sur deux sessions de validation de vingt minutes~; le modèle de chaque horizon est calibré sur quatre autres sessions. Données~:
\texttt{hf\_alpha.ic\_by\_horizon}.}\label{fig:hf:short-horizon-alpha:ic}
\end{omfigure}

Une prévision est aussi une course contre son propre âge. La même prévision à deux secondes, construite à partir de variables âgées de $k$ messages, conserve un
coefficient d'information de 0.477 à un message de retard (une médiane de 33 millisecondes ici), 0.405 à vingt messages de retard (0.9 seconde) et 0.303 à cinquante messages de retard
(2.3 secondes). À l'horizon de deux secondes, la décroissance est lente~; ce qui décroît vite, c'est la valeur d'agir avant les autres, qui fait l'objet du chapitre~9.

\section{Lier la prévision aux cotations~: décalage}

\begin{definition}[Décalage de prévision]\label{def:hf:short-horizon-alpha:skew}
Un \emph{décalage de prévision}\index{décalage de prévision} (forecast skew) déplace les cotations d'un teneur de marché dans le sens de sa prévision à court horizon~: le côté que la
prévision désigne comme devant être pris pour cible (le cours vendeur lorsque le prix doit monter) est écarté du \omterm{def:hf:fair-value:fair}{prix juste} ou retiré, le côté opposé étant
amélioré.
\end{definition}

Dans un marché à petit tick, le décalage est continu~: la dérive de la prévision entre dans le problème de contrôle (chapitre~4, \omterm{def:hf:extensions-of-the-inventory-framework:drift}{cotation ajustée de la dérive}), ou le \omterm{def:hf:fair-value:fair}{prix juste} lui-même
se déplace (chapitre~2). Dans un marché à un seul tick, c'est une décision sur le côté à afficher. Cartea, Donnelly et Jaimungal ont introduit un signal de déséquilibre de volume
dans le problème de contrôle d'un teneur de marché, l'ont calibré sur onze actions du Nasdaq sur un semestre, l'ont testé sur le semestre suivant et ont constaté que
le signal a amélioré la stratégie hors échantillon.

\begin{omfigure}
\begin{tikzpicture}[font=\small,>=stealth]
\draw[->] (-6,0) -- (6.3,0);
\node[font=\footnotesize] at (0,-1.15) {prévision $\hat\alpha$ (ticks)};
\foreach \x/\l in {-4.5/{$-0.9$},-1.5/{$-0.3$},0/{0},1.5/{$0.3$},4.5/{$0.9$}} {\draw (\x,0.1) -- (\x,-0.1) node[below=1pt,font=\footnotesize] {\l};}
\fill[omDef!15] (-1.5,0.25) rectangle (1.5,1.3);
\node[font=\footnotesize] at (0,0.77) {les deux côtés};
\fill[omThm!15] (1.5,0.25) rectangle (4.5,1.3);
\node[font=\footnotesize] at (3,0.77) {le bid seul};
\fill[omThm!30] (4.5,0.25) rectangle (6.1,1.3);
\node[font=\footnotesize,align=center] at (5.3,0.77) {le bid,\\ et achat};
\fill[omProp!15] (-4.5,0.25) rectangle (-1.5,1.3);
\node[font=\footnotesize] at (-3,0.77) {l'ask seul};
\fill[omProp!30] (-6,0.25) rectangle (-4.5,1.3);
\node[font=\footnotesize,align=center] at (-5.25,0.77) {l'ask,\\ et vente};
\node[font=\scriptsize,anchor=north] at (3,-0.55) {décalage};
\node[font=\scriptsize,anchor=north] at (5.3,-0.55) {seuil d'agression};
\end{tikzpicture}

\omcaption{Le couplage du chapitre entre une prévision à deux secondes et des cotations d'un tick~: les deux côtés sont affichés tant que la prévision reste dans $\pm0.3$ ticks~;
au-delà, le côté qui serait exposé à une sélection adverse est retiré (skew)~; au-delà de $\pm0.9$ ticks, la moitié de la fourchette, plus les frais de preneur et une marge, le
teneur de marché traverse aussi la fourchette pour un lot (\omterm{def:hf:short-horizon-alpha:take}{seuil d'agression}).}\label{fig:hf:short-horizon-alpha:coupling}
\end{omfigure}

\section{Quand agresser plutôt que coter}

\begin{definition}[Seuil d'agression]\label{def:hf:short-horizon-alpha:take}
Un \emph{seuil d'agression}\index{seuil d'agression} (take threshold) est la taille de la prévision à court horizon au-delà de laquelle un teneur de marché traverse la fourchette au lieu de
attendre une exécution~: au moins la moitié de la fourchette, plus les frais de preneur, plus une marge pour l'erreur de la prévision et pour les autres traders qui voient le
même signal.
\end{definition}

Un teneur de marché qui achète au cours vendeur paie la moitié de la fourchette et les frais de preneur par rapport au mid~; il ne gagne que si le mid monte davantage avant que la
prévision ne soit périmée. Ici, la moitié de la fourchette vaut un demi-tick, les frais de preneur un dixième de tick, la marge trois dixièmes~: 0.9 ticks. Des prévisions aussi grandes sont rares~:
sur les séances de validation, 0.01~\% des prévisions dépassent ce seuil en valeur absolue, contre 5.8~\% pour le seuil de décalage de 0.3.

Sur six séances de test, un teneur de marché à un lot en B (frais de preneur de 0.1 cent par action, sans rebate)~:

\begin{center}
\footnotesize
\setlength{\tabcolsep}{4pt}%
\begin{tabular}{@{}>{\raggedright\arraybackslash}p{3.4cm}>{\centering\arraybackslash}p{2.1cm}>{\centering\arraybackslash}p{2.1cm}>{\centering\arraybackslash}p{2.1cm}>{\centering\arraybackslash}p{2.1cm}>{\centering\arraybackslash}p{2.1cm}@{}}
\toprule
couplage & actions passives & mark-out passif & agressions & mark-out des agressions & P\&L (écart-type)\\
& par séance & (ticks par action) & par séance & (ticks par action) & (\$ par séance)\\
\midrule
aucun & 16\,550 & 0.294 & & & $-33.58$ (107.68)\\
skew & 11\,050 & 0.451 & & & 25.58 (50.82)\\
skew et agression & 10\,867 & 0.452 & 4.2 & 0.540 & 46.17 (51.22)\\
idem, un message en retard & 10\,883 & 0.414 & 4.5 & 0.722 & 46.47 (36.29)\\
\bottomrule
\end{tabular}
\end{center}

Le skew relève le mark-out passif de 0.16 ticks par action et abandonne un tiers du volume, le même compromis que le fading du chapitre~6, à partir d'une
prévision de direction plutôt que d'une prévision de mark-out. Les agressions sont rares, quatre par séance~; leur mark-out est positif mais repose sur 25 à 27 transactions, et
le chiffre plus élevé de la version tardive relève du bruit. Avec un message de retard, le mark-out passif baisse de 8~\%. La colonne P\&L est trop bruitée sur six séances pour classer
les couplages, comme dans les chapitres précédents.

\section{Fiches de stratégie}

\begin{strategyfile}{Skew de cotation sur le déséquilibre du carnet}\label{strat:hf:short-horizon-alpha:imbalance}
\sfield{Qui paie, et pourquoi} Des preneurs non informés qui traversent vers le côté du carnet qui va tenir~; les preneurs informés sont évités au lieu
d'être rémunérés.
\sfield{Instruments et plateformes} Actions à grand tick, contrats à terme et fonds indiciels cotés (ETF) sur des plateformes à priorité prix-temps.
\sfield{Signal} Déséquilibre des files $(Q^b-Q^a)/(Q^b+Q^a)$ aux meilleurs prix, avec le déséquilibre du flux d'ordres (OFI) sur une courte fenêtre~; prévision $\hat\alpha$ de la
variation du mid sur une à quelques secondes.
\sfield{Dimensionnement et exécution} Afficher les deux côtés tant que $|\hat\alpha|$ est faible~; retirer le côté opposé à la prévision au-delà d'un seuil~; hystérésis pour
conserver ses positions dans la file (chapitre~6).
\sfield{Coûts} Positions dans la file perdues à chaque retrait~; messages~; le volume abandonné.
\sfield{Comment elle s'éteint} Chaque teneur de marché lit le même déséquilibre, si bien que le côté appelé à tenir est encombré et que la file qui subsiste est celle qui va
être consommée~; aucun épisode public précis n'est documenté.
\sfield{Horizon, capacité, infrastructure} Secondes~; un à quelques lots par niveau~; un flux ordre par ordre et une mise à jour rapide des cotations.
\sfield{Backtester honnêtement} Un modèle de position dans la file (livre~7, chapitre~18) ou un simulateur réactif~: des ordres fantômes qui ne perdent jamais leur place
rendent le retrait trop favorable.
\sfield{Sources} Cartea, Donnelly et Jaimungal (2018)~; Kolm, Turiel et Westray (2023)~; ce chapitre~: mark-out passif 0.451 contre 0.294 ticks par
action, sur un marché synthétique (borne supérieure).
\end{strategyfile}

\begin{strategyfile}{Agression sur déséquilibre du flux d'ordres}\label{strat:hf:short-horizon-alpha:ofi}
\sfield{Qui paie, et pourquoi} Des fournisseurs de liquidité dont les cotations sont périmées après une rafale de flux d'ordres.
\sfield{Instruments et plateformes} Contrats à terme et actions liquides, avec des carnets profonds.
\sfield{Signal} Déséquilibre du flux d'ordres (OFI) sur la dernière seconde, normalisé par la profondeur (relation linéaire de Cont, Kukanov et Stoikov).
\sfield{Dimensionnement et exécution} Traverser la fourchette pour un lot lorsque $|\hat\alpha|$ dépasse le \omterm{def:hf:short-horizon-alpha:take}{seuil d'agression}~; un délai de repos entre deux agressions.
\sfield{Coûts} La moitié de la fourchette et les frais de preneur sur chaque transaction, à comparer au gain de la prévision.
\sfield{Comment elle s'éteint} Le même signal est visible de chaque trader rapide~; le premier à agir saisit l'opportunité, et la course aux armements raccourcit
sa durée de vie (chapitre~9).
\sfield{Horizon, capacité, infrastructure} De la sous-seconde à quelques secondes~; capacité très faible par opportunité~; colocation.
\sfield{Backtester honnêtement} Exécuter un message après le signal, au prix alors disponible~; ne compter que les agressions dont la prévision a dépassé le
seuil à l'instant différé.
\sfield{Sources} Cont, Kukanov et Stoikov (2014)~; ce chapitre~: 4.2 agressions par séance, mark-out 0.540 ticks sur trop peu de transactions pour s'y fier.
\end{strategyfile}

\begin{strategyfile}{Continuation du signe de la transaction}\label{strat:hf:short-horizon-alpha:sign}
\sfield{Qui paie, et pourquoi} Le teneur de marché évite le côté que prend un métaordre~; le propriétaire du métaordre paie par la fourchette qu'il ne cesse
de payer.
\sfield{Instruments et plateformes} Tout instrument à mémoire longue dans les signes des ordres.
\sfield{Signal} Volume agressif net sur une fenêtre~; l'autocorrélation des signes du livre~7, chapitre~9.
\sfield{Dimensionnement et exécution} Un terme dans la prévision~; pris seul, la plus faible des cinq variables retenues ici (coefficient d'information 0.077).
\sfield{Coûts} Comme pour le skew.
\sfield{Comment elle s'éteint} Les algorithmes d'exécution qui randomisent les tailles et le calendrier masquent le signe~; la persistance est en partie compensée par des ordres à cours limité
à retour à la moyenne (Bouchaud et coauteurs).
\sfield{Horizon, capacité, infrastructure} De quelques secondes à quelques minutes.
\sfield{Backtester honnêtement} Attribuer leur signe aux transactions avec la règle que le système en production utilisera (livre~7, chapitre~9), et non avec la vérité du simulateur.
\sfield{Sources} Bouchaud, Gefen, Potters et Wyart (2004)~; le coefficient de ce chapitre.
\end{strategyfile}

\begin{strategyfile}{Cotation guidée par les contrats à terme}\label{strat:hf:short-horizon-alpha:futures}
\sfield{Qui paie, et pourquoi} Des fournisseurs de liquidité plus lents dans le fonds ou sur les actions, dont les cotations sont en retard sur le contrat à terme.
\sfield{Instruments et plateformes} Un contrat à terme sur indice et son fonds ou ses composantes~; un contrat à terme sur Treasury et l'obligation au comptant~; tout meneur et tout suiveur.
\sfield{Signal} La variation du mid du meneur sur la dernière fenêtre, traduite dans les unités du suiveur (\omterm{def:hf:fair-value:cross}{prix juste inter-instruments} du chapitre~2).
\sfield{Dimensionnement et exécution} Retirer la cotation du suiveur du côté contre lequel le meneur a bougé~; agresser lorsque l'écart dépasse le \omterm{def:hf:short-horizon-alpha:take}{seuil d'agression}.
\sfield{Coûts} La latence entre les deux plateformes~; les frais sur les deux.
\sfield{Comment elle s'éteint} Vitesse~: la durée de vie médiane d'un arbitrage entre contrat à terme sur indice et fonds est passée de 97 millisecondes en 2005 à 7 en 2011, à mesure que les firmes investissaient dans
la vitesse (Budish, Cramton et Shim), tandis que son profit médian par opportunité est resté proche de 0.08 points d'indice.
\sfield{Horizon, capacité, infrastructure} De la milliseconde à la seconde~; un chemin de données depuis la plateforme du meneur (chapitre~8).
\sfield{Backtester honnêtement} Horodatages des deux plateformes sur une même horloge, avec le délai de propagation réel entre elles.
\sfield{Sources} Budish, Cramton et Shim (2015)~; ce chapitre~: le mouvement du meneur a un coefficient d'information de 0.288 à deux secondes.
\end{strategyfile}

\begin{strategyfile}{Retour à la moyenne après balayage}\label{strat:hf:short-horizon-alpha:sweep}
\sfield{Qui paie, et pourquoi} Un trader qui a balayé plusieurs niveaux et poussé le prix au-delà de son niveau, lorsque l'impact se dissipe.
\sfield{Instruments et plateformes} Actions et contrats à terme à impact transitoire.
\sfield{Signal} La variation propre de l'instrument sur la dernière fenêtre, après un balayage~; le propagateur du livre~10, chapitre~12 donne le
retour attendu à la moyenne.
\sfield{Dimensionnement et exécution} Coter dans le sens du retour~: afficher le côté qui s'oppose au balayage.
\sfield{Coûts} Sélection adverse lorsque le balayage était informé.
\sfield{Comment elle s'éteint} Lorsque le balayage était une information et non de la liquidité~: le prix poursuit son mouvement. Sur le marché de ce chapitre, la variation propre prédit
une poursuite du mouvement (coefficient $+0.27$), il n'y a donc pas de retour à exploiter.
\sfield{Horizon, capacité, infrastructure} Secondes.
\sfield{Backtester honnêtement} Séparer les balayages selon ce qui les a suivis, uniquement hors échantillon (out-of-sample)~; le marché artificiel ici montrerait la stratégie perdante.
\sfield{Sources} Bouchaud, Gefen, Potters et Wyart (2004)~; le livre~10, chapitre~12.
\end{strategyfile}

\section{Tutoriel~: prévoir le fonds à partir du futur}\label{tut:hf:short-horizon-alpha:tut}

\begin{tutorial}
\textbf{Objectif.} Construire une prévision à quelques secondes à partir du carnet, des flux et d'un instrument précurseur, l'évaluer hors échantillon et la coupler à un teneur de cotations.
\textbf{État final~:} les trois tableaux, \cref{fig:hf:short-horizon-alpha:ic,fig:hf:short-horizon-alpha:calib}.
\begin{tutsteps}
\item \textbf{Variables.} \texttt{firm.hfalpha.FeatureEngine} se met à jour à chaque message~; le même code s'exécute sur les sessions enregistrées (\texttt{stream})
et à l'intérieur du teneur de cotations.
\omcode{code/firm/hfalpha/firm_hfalpha.py}{56}{79}{Déséquilibre du flux d'ordres, flux de transactions, mouvement propre et déséquilibre, mis à jour à chaque message.}\label{lst:hf:short-horizon-alpha:features}
\item \textbf{Ajuster et évaluer.} \texttt{hf\_alpha.model(h)} ajuste une régression ridge sur quatre sessions~; \texttt{ics(h)} la note sur deux autres~;
\texttt{ic\_by\_delay()} vieillit les variables.
\item \textbf{Coupler.} \texttt{AlphaQuoter} lit le carnet sans ses propres ordres, prévoit, décale et agresse.
\omcode{code/firm/hfalpha/firm_hfalpha.py}{145}{166}{La prévision détermine le côté à afficher et le moment où agresser.}\label{lst:hf:short-horizon-alpha:quoter}
\item \textbf{Comparer} avec \texttt{compare()} sur six sessions de test.
\end{tutsteps}
\textbf{Pour aller plus loin.} Remplacer la régression ridge par un boosting de gradient (livre~12, chapitre~5) et vérifier que le coefficient hors échantillon
s'améliore~; rendre le retard du précurseur aléatoire d'une session à l'autre et suivre le coefficient du précurseur.
\end{tutorial}

\begin{omfigure}
\begin{tikzpicture}
\begin{axis}[width=11cm,height=5.2cm,xlabel={\small dixième des prévisions},ylabel={\small ticks sur deux secondes},xtick={1,2,3,4,5,6,7,8,9,10},
  xmin=0.5,xmax=10.5,ymin=-0.4,ymax=0.4,tick label style={font=\footnotesize},grid=major,grid style={gray!25},table/col sep=comma,
  legend style={font=\scriptsize,at={(0.5,-0.3)},anchor=north},legend columns=2]
\addplot[omDef,thick,mark=*,mark size=1.2pt] table[x=group,y=pred]{figdata/hft/07-short-horizon-alpha/calibration.csv};
\addplot[omThm,thick,only marks,mark=square*,mark size=1.5pt] table[x=group,y=real]{figdata/hft/07-short-horizon-alpha/calibration.csv};
\legend{prévision moyenne,variation réalisée moyenne}
\end{axis}
\end{tikzpicture}

\omcaption{La prévision à deux secondes sur les sessions de validation, classées par dixièmes~: prévision moyenne et variation réalisée moyenne du mid du fonds,
en ticks. Données~: \texttt{hf\_alpha.calibration}.}\label{fig:hf:short-horizon-alpha:calib}
\end{omfigure}

\section{Construction~: le moteur d'alpha à court horizon}\label{bld:hf:short-horizon-alpha:hfalpha}

\begin{build}
\bfield{Rôle} Un seul moteur de variables pour la recherche et le trading, une prévision ajustée et évaluée hors échantillon, et les couplages qu'utilise un teneur de marché.
\bfield{Interface} \texttt{FeatureEngine(window, lot, leader).update(t, kind, agg, qty, bid, bid\_qty, ask, ask\_qty)}~;
\texttt{stream(tape, window, leader)}~; \texttt{target(times, mid, horizon)}~; \texttt{Ridge.fit(X, y, lam)}, \texttt{.predict}~; \texttt{ic}~;
\texttt{AlphaQuoter(model, window, leader, skew, take, lag, size, limit)}.
\bfield{Règles} Le moteur ne voit que ce qu'un flux a montré~; le précurseur est lu à la même date. Les variables du teneur de cotations utilisent le carnet sans ses
propres ordres, comme le modèle a été ajusté. Chaque score est hors échantillon.
\bfield{Tests d'acceptation} \texttt{code/firm/hfalpha/tests/}~: déséquilibre du flux d'ordres, flux de transactions, mouvement du précurseur et mouvement propre à la main sur un flux
de cinq messages, y compris le vidage de la fenêtre~; la cible à la main~; la régression ridge retrouve une relation artificielle.
\bfield{Pour aller plus loin} Variables sur plusieurs niveaux du carnet~; une prévision par horizon, mélangée selon la durée de détention du teneur de cotations.
\end{build}

\begin{omsources}
\item R.~Cont, A.~Kukanov et S.~Stoikov, ``The price impact of order book events'', \emph{Journal of Financial Econometrics} 12(1), 2014.
\item P.~N.~Kolm, J.~Turiel et N.~Westray, ``Deep order flow imbalance: extracting alpha at multiple horizons from the limit order book'',
\emph{Mathematical Finance} 33(4), 2023.
\item Á.~Cartea, R.~Donnelly et S.~Jaimungal, ``Enhancing trading strategies with order book signals'', \emph{Applied Mathematical Finance} 25(1),
2018.
\item E.~Budish, P.~Cramton et J.~Shim, ``The high-frequency trading arms race: frequent batch auctions as a market design response'',
\emph{Quarterly Journal of Economics} 130(4), 2015.
\item J.-P.~Bouchaud, Y.~Gefen, M.~Potters et M.~Wyart, ``Fluctuations and response in financial markets: the subtle nature of random price
changes'', \emph{Quantitative Finance} 4(2), 2004.
\end{omsources}

\section{Exercices}

\begin{exercise}[$\star$]\label{exo:hf:short-horizon-alpha:1}
La file acheteuse compte 12 lots et la file vendeuse 4. Avec le coefficient du chapitre de 0.30 tick par unité de déséquilibre, que
prévoit le seul déséquilibre~?
\end{exercise}

\begin{exercise}[$\star$]\label{exo:hf:short-horizon-alpha:2}
Avec un tick d'un centime, une demi-fourchette d'un demi-tick et des frais de preneur de liquidité de 0.1 centime par action, quelle est la plus petite prévision pour laquelle
l'agression est tout juste rentable, avant toute marge~?
\end{exercise}

\begin{exercise}[$\star$]\label{exo:hf:short-horizon-alpha:3}
Le meilleur cours acheteur passe de 99 à 100 avec 3 lots, tandis que le meilleur cours vendeur reste à 101 et passe de 5 à 6 lots. Quel est l'accroissement du déséquilibre du flux d'ordres~?
\end{exercise}

\begin{exercise}[$\star\star$]\label{exo:hf:short-horizon-alpha:4}
Pourquoi la variable de mouvement propre a-t-elle un coefficient positif sur ce marché, et que faudrait-il attendre sur une action réelle~?
\end{exercise}

\begin{exercise}[$\star\star$]\label{exo:hf:short-horizon-alpha:5}
D'après la \cref{fig:hf:short-horizon-alpha:ic}, quelle variable est la meilleure à une demi-seconde et laquelle à dix secondes~? Que suggère cela pour le
couplage d'un teneur de cotations qui détient ses positions pendant cinq secondes~?
\end{exercise}

\begin{exercise}[$\star\star$]\label{exo:hf:short-horizon-alpha:6}
Pourquoi une prévision dont le coefficient d'information vaut 0.48 peut-elle rarement payer le franchissement de la fourchette~?
\end{exercise}

\begin{exercise}[$\star\star\star$]\label{exo:hf:short-horizon-alpha:7}
\emph{Programmation.} Avec \texttt{ic\_by\_delay()}, donner le coefficient d'information de la prévision à deux secondes faite dix messages en retard, et le temps médian
que mettent ces messages~?
\end{exercise}

\begin{exercise}[$\star\star\star$]\label{exo:hf:short-horizon-alpha:8}
\emph{Trouver la faille.} «~Notre backtest agresse dès que la prévision dépasse la fourchette et gagne 0.7 tick par transaction~; il exécute au mid au
moment de la prévision.~»
\end{exercise}

\section{Problème~: trois contre un sur le bid}

\begin{problem}[{Problème du week-end --- trois contre un sur le bid}]\label{pb:hf:short-horizon-alpha:1}
Le teneur de marché d'un fonds doit décider, message par message, laquelle de ses deux cotations afficher et quand agresser.

\medskip\noindent\textbf{Partie I --- La prévision.}
\begin{enumerate}
\item Définir l'\omterm{def:hf:short-horizon-alpha:alpha}{alpha à court horizon} et expliquer en quoi il diffère de l'alpha intrajournalier.
\item Lister les cinq variables et la manière dont chacune est calculée.
\item Donner les coefficients d'information hors échantillon à deux secondes, ainsi que celui de la combinaison.
\item Interpréter les coefficients par unité du déséquilibre et du mouvement du précurseur.
\item Pourquoi ces coefficients constituent-ils une borne supérieure pour un marché réel~?
\end{enumerate}

\medskip\noindent\textbf{Partie II --- Horizons et ancienneté.}
\begin{enumerate}[resume]
\item Comment le coefficient de la combinaison évolue-t-il d'une demi-seconde à dix secondes~?
\item Qu'ont établi Kolm, Turiel et Westray au sujet de l'horizon effectif~?
\item Donner le coefficient pour un, vingt et cinquante messages de retard, et les durées correspondantes.
\item Que montre la \cref{fig:hf:short-horizon-alpha:calib}~?
\end{enumerate}

\medskip\noindent\textbf{Partie III --- Utilisation.}
\begin{enumerate}[resume]
\item Définir un \omterm{def:hf:short-horizon-alpha:skew}{décalage de prévision} et le décrire sur un marché à un tick.
\item Définir le \omterm{def:hf:short-horizon-alpha:take}{seuil d'agression} et en déduire les 0.9~tick du chapitre.
\item À quelle fréquence la prévision dépasse-t-elle chaque seuil~?
\item Donner les parts passives et les mark-outs des couplages.
\end{enumerate}

\medskip\noindent\textbf{Partie IV --- Le verdict.}
\begin{enumerate}[resume]
\item Énoncer le \emph{résultat nommé}~: le mark-out passif par action du teneur de marché décalé, comparé à celui du teneur non décalé, et la fraction de celui-ci perdue
un message plus tard.
\item Pourquoi les résultats des agressions ne sont-ils pas une preuve~?
\item Quel est le lien entre le décalage et l'estompement du chapitre~6~?
\item Quelle fiche de stratégie n'a aucun edge sur ce marché, et pourquoi~?
\item Comment la course à l'armement a-t-elle modifié le trading mené par les contrats à terme entre 2005 et 2011~?
\item Que changeriez-vous dans le backtest avant de croire au mark-out des agressions~?
\item En une phrase~: à quoi sert surtout une prévision de quelques secondes pour un teneur de marché~?
\end{enumerate}
\end{problem}

\section{Questions d'entretien}

\begin{interviewq}[$\star$ \iqroles{trader}]\label{iq:hf:short-horizon-alpha:1}
Le contrat à terme vient de monter d'un tick et votre cotation du fonds n'a pas bougé. Que faites-vous de votre cours vendeur, et pourquoi~?
\end{interviewq}

\begin{interviewq}[$\star\star$ \iqroles{chercheur}]\label{iq:hf:short-horizon-alpha:2}
Définissez le déséquilibre du flux d'ordres aux meilleurs prix et expliquez pourquoi il prédit mieux les variations de prix que le volume signé des transactions.
\end{interviewq}

\begin{interviewq}[$\star\star$ \iqroles{chercheur}]\label{iq:hf:short-horizon-alpha:3}
Comment évaluez-vous une prévision dont les fenêtres cibles se chevauchent, message après message~?
\end{interviewq}

\begin{interviewq}[$\star\star$ \iqroles{développeur}]\label{iq:hf:short-horizon-alpha:4}
La recherche calcule les variables à partir des données stockées, la production les calcule à partir du flux en direct. Comment vous assurez-vous que ce sont les mêmes valeurs~?
\end{interviewq}

\begin{interviewq}[$\star\star$ \iqroles{trader, chercheur}]\label{iq:hf:short-horizon-alpha:5}
Quand un teneur de marché doit-il agresser plutôt qu'afficher~? Donnez le seuil et ce qui entre dans sa marge.
\end{interviewq}

\begin{interviewq}[$\star\star\star$ \iqroles{chercheur}]\label{iq:hf:short-horizon-alpha:6}
Une prévision $\hat\alpha$ d'une variation de prix $y$ a pour corrélation $\rho$ avec elle, et $\hat\alpha$ est l'espérance conditionnelle. Si $y$ a un écart-type
$s$, à quelle fréquence, approximativement, $|\hat\alpha|$ dépasse-t-il un seuil $c$ lorsque $\hat\alpha$ est gaussienne~? Appliquez cela à $\rho=0.48$, $s=0.6$~tick, $c=0.9$.
\end{interviewq}
